\documentclass[11pt]{article}

\usepackage[margin=1in]{geometry}
\usepackage{booktabs}
\usepackage{amsmath}
\usepackage{hyperref}

\title{Dual-Domain MMD Adaptation: Regularization Ablation in the Large-Source,
Small-Target Regime}
\author{}
\date{\today}

\begin{document}
\maketitle

\section{Setup}

All five variants below share the same starting checkpoint (syn\_large pretrain
baseline) and the same target set (real\_small), differing only in whether/how MMD
regularization was applied:

\begin{itemize}
    \item \textbf{real\_small\_from\_syn\_large} (no-MMD control): plain continued
    fine-tune, no MMD term, no source batches. 100 epochs, batch 8.
    \item \textbf{syn\_large\_source\_real\_small\_target} (unregularized MMD):
    constant mmd\_weight = 0.8 (default), no bandwidth freeze. 50 epochs, batch 8. See
    \texttt{experiment\_summary\_v4.tex}.
    \item \textbf{syn\_large\_source\_real\_small\_target\_reg}: mmd\_weight decaying
    linearly 1.0 $\to$ 0.1, bandwidth freeze at epoch 1 (recalled from memory, not
    logged with certainty -- may have been epoch 2). 50 epochs.
    \item \textbf{\_reg\_smaller\_weight}: mmd\_weight decaying linearly 0.9 $\to$
    0.02, bandwidth freeze at epoch 1. 50 epochs, batch 4.
    \item \textbf{\_reg\_smaller\_weight\_longer\_run\_final}: same weight schedule as
    \_reg\_smaller\_weight (0.9 $\to$ 0.02, freeze at epoch 1), but 100 epochs instead
    of 50.
\end{itemize}

\paragraph{A note on hyperparameter provenance.} None of the MMD-specific settings
(weight, schedule, freeze epoch) are captured in ultralytics' own \texttt{args.yaml} --
they are popped from the overrides dict before being handed to the stock trainer, so
they leave no automatic record. The values above for the three regularized runs are
as recalled after the fact rather than read from a log; treat the base \_reg run's
freeze epoch in particular as approximate.

\section{Results}

\begin{table}[h!]
    \centering
    \begin{tabular}{lrrrr}
        \toprule
        Model & person mAP50 & person mAP50-95 & vehicle mAP50 & vehicle mAP50-95 \\
        \midrule
        real\_small\_from\_syn\_large (no-MMD) & 0.644 & 0.303 & 0.877 & 0.626 \\
        syn\_large\_source\_real\_small\_target (unreg.\ MMD) & 0.642 & 0.296 & 0.877 & 0.625 \\
        \_reg & 0.649 & 0.300 & 0.881 & 0.628 \\
        \_reg\_smaller\_weight & 0.646 & 0.302 & 0.879 & 0.625 \\
        \_reg\_smaller\_weight\_longer\_run\_final & \textbf{0.651} & \textbf{0.304} & \textbf{0.882} & \textbf{0.630} \\
        \bottomrule
    \end{tabular}
    \caption{Validated on the real val set -- the target domain for this direction,
    the metric of primary interest.}
\end{table}

\begin{table}[h!]
    \centering
    \begin{tabular}{lrrrr}
        \toprule
        Model & person mAP50 & person mAP50-95 & vehicle mAP50 & vehicle mAP50-95 \\
        \midrule
        real\_small\_from\_syn\_large (no-MMD) & 0.129 & 0.0349 & 0.222 & 0.146 \\
        syn\_large\_source\_real\_small\_target (unreg.\ MMD) & 0.111 & 0.0266 & 0.214 & 0.138 \\
        \_reg & 0.103 & 0.0255 & 0.218 & 0.141 \\
        \_reg\_smaller\_weight & 0.116 & 0.0282 & 0.227 & 0.145 \\
        \_reg\_smaller\_weight\_longer\_run\_final & 0.0886 & 0.0216 & 0.215 & 0.140 \\
        \bottomrule
    \end{tabular}
    \caption{Validated on the syn val set -- the source (forgotten) domain for this
    direction.}
\end{table}

\section{Discussion}

\textbf{Regularization recovers a small, consistent MMD benefit on the target
domain.} In \texttt{experiment\_summary\_v4.tex}, unregularized MMD showed no visible
advantage over the no-MMD control on vehicle/real (0.877 vs.\ 0.877 mAP50, identical).
Here, all three regularized variants edge out both baselines on this exact metric:
0.881, 0.879, and 0.882 vs.\ 0.877 for the two baselines. This is a small margin
(0.2--0.6 points), but it is now three independent confirmations pointing the same
direction rather than one ambiguous run -- a materially stronger signal than the
unregularized case gave. The best performer overall is
\_reg\_smaller\_weight\_longer\_run\_final (100 epochs, smaller decaying weight),
which leads on both vehicle \emph{and} person, on both val sets simultaneously.

\textbf{The source (forgetting) domain shows no consistent pattern.} On syn,
\_reg\_smaller\_weight has the best vehicle score of all five variants (0.227), but
\_reg\_smaller\_weight\_longer\_run\_final -- the best performer on the target domain
-- has one of the weaker person scores on syn (0.0886, worse than every other
variant). Regularization does not appear to uniformly help forgetting the way it
helps the target-domain result; the effect there is mixed rather than directional.

\textbf{Caveats.} Margins throughout are small (a few tenths of a point), these are
single runs with no repeated seeds, and the exact hyperparameters for the base \_reg
run are recalled rather than logged (see note in Setup). The directional consistency
across three separately-configured regularized runs is encouraging, but confirming
this with a repeat run or a small sweep around the smaller-weight/longer-run setting
would be the natural next step before treating the magnitude of the effect as settled.

\end{document}
